\documentclass[
    language=english,
    doctype=manual,
    institution=none,
    titlestyle=book,
    oneside,
]{../../omnilatex}

\RequirePackage{config/document-settings}

\addbibresource{../../bib/bibliography.bib}

\begin{document}

\maketitle

\tableofcontents

%=============================================================================
\chapter{Getting Started}
\label{sec:getting-started}
%=============================================================================

This guide is for contributors to the DataPipe project. It covers the
repository layout, development environment setup, coding standards, and
the pull request workflow. Before contributing, please read the Code of
Conduct and the existing issues to find a task that matches your
experience level.

\section{Development Environment}

\subsection{Prerequisites}

\begin{itemize}
    \item JDK 21 (Eclipse Temurin recommended).
    \item Docker Desktop 4.x or Docker Engine 24+.
    \item \texttt{gh} CLI for GitHub interaction.
    \item \texttt{pre-commit} for commit hooks.
    \item \texttt{just} as a command runner (alternative to Make).
\end{itemize}

\subsection{Cloning and Building}

\begin{lstlisting}[language=bash, caption={Repository setup and first build.}]
git clone https://github.com/datapipe/datapipe.git
cd datapipe
./gradlew build          # Full build with tests
./gradlew assemble       # Compile only (faster)
./gradlew test           # Run unit tests
\end{lstlisting}

\subsection{IDE Setup}

Import the project as a Gradle project. IntelliJ IDEA 2024+ and VS Code
with the Java Extension Pack are officially supported. Enable the
following settings:

\begin{lstlisting}[language=bash, caption={Recommended IDE settings.}]
# IntelliJ: enable annotation processing
# Settings > Build > Compiler > Annotation Processors > Enable

# VS Code: install extensions
code --install-extension vscjava.vscode-java-pack
code --install-extension vscjava.vscode-java-debug
\end{lstlisting}

%=============================================================================
\chapter{Architecture Overview}
\label{sec:architecture}
%=============================================================================

\section{High-Level Components}

DataPipe is organised into four major subsystems. The following diagram
shows their relationships:

\begin{figure}[htbp]
    \centering
    \begin{tikzpicture}[
        node distance=1.4cm and 2cm,
        box/.style={draw, rounded corners, minimum width=3cm,
                     minimum height=1cm, text centered, font=\small,
                     fill=#1!20, draw=#1!60!black},
        arr/.style={->, >=stealth, thick},
    ]
        \node[box=blue]   (cli)   {CLI};
        \node[box=green]  (core)  [below=of cli] {Core Runtime};
        \node[box=orange] (conn)  [below left=of core] {Connectors};
        \node[box=purple] (trans) [below right=of core] {Transforms};
        \node[box=red]    (admin) [right=of core] {Admin API};

        \draw[arr] (cli) -- (core);
        \draw[arr] (core) -- (conn);
        \draw[arr] (core) -- (trans);
        \draw[arr] (admin) -- (core);
    \end{tikzpicture}
    \caption{DataPipe subsystem architecture.}
    \label{fig:subsystem-arch}
\end{figure}

\section{Module Structure}

The repository is a multi-module Gradle project:

\begin{lstlisting}[language=, caption={Top-level directory layout.}]
datapipe/
  core/           # Core runtime and pipeline engine
  connectors/     # Source and sink connector plugins
  transforms/     # Built-in transform library
  cli/            # Command-line interface
  admin-api/      # REST admin API
  web-dashboard/  # React frontend
  common/         # Shared utilities and types
  docs/           # Documentation source files
  deploy/         # Helm charts and Docker configs
\end{lstlisting}

\section{Data Flow}

A record passes through the following stages:

\begin{enumerate}
    \item \textbf{Ingestion:} A connector reads a record from the source.
    \item \textbf{Deserialisation:} The record is decoded from bytes into
        an internal \texttt{Record} object.
    \item \textbf{Transformation:} Zero or more transforms are applied.
    \item \textbf{Serialisation:} The record is encoded for the sink.
    \item \textbf{Delivery:} The connector writes the record to the
        destination.
\end{enumerate}

%=============================================================================
\chapter{Coding Standards}
\label{sec:coding-standards}
%=============================================================================

\section{Java Code Style}

DataPipe follows the Google Java Style Guide with the following
modifications:

\begin{itemize}
    \item Maximum line length: 100 characters.
    \item Four-space indentation (no tabs).
    \item \texttt{@Override} annotations are mandatory on all
        overriding methods.
    \item All public APIs require Javadoc.
\end{itemize}

\section{Commit Messages}

Follow the Conventional Commits specification:

\begin{lstlisting}[language=bash, caption={Commit message examples.}]
feat(connector): add Snowflake sink connector
fix(core): handle null records in window aggregator
docs: update contributing guide with PR checklist
test(transform): add edge cases for currency converter
\end{lstlisting}

\section{Test Requirements}

\begin{table}[htbp]
    \centering
    \caption{Test coverage requirements by module.}
    \label{tab:test-coverage}
    \begin{tabular}{@{}lrr@{}}
        \toprule
        Module & Minimum Coverage & Critical Path \\
        \midrule
        core/ & 90\% & Yes \\
        connectors/ & 85\% & Yes \\
        transforms/ & 80\% & Yes \\
        cli/ & 75\% & No \\
        admin-api/ & 85\% & Yes \\
        \bottomrule
    \end{tabular}
\end{table}

%=============================================================================
\chapter{Testing}
\label{sec:testing}
%=============================================================================

\section{Unit Tests}

Unit tests live alongside source files in \texttt{src/test/java}. Use
JUnit 5 and Mockito:

\begin{lstlisting}[language=java, caption={Example unit test.}]
@Test
void shouldFilterRecordsByPredicate() {
    FilterTransform transform = new FilterTransform(
        Config.of("predicate", "status == 'active'")
    );

    Record active = Record.of("status", "active");
    Record inactive = Record.of("status", "inactive");

    assertThat(transform.apply(active)).isPresent();
    assertThat(transform.apply(inactive)).isEmpty();
}
\end{lstlisting}

\section{Integration Tests}

Integration tests require Docker and run in CI. They verify end-to-end
pipeline execution against real Kafka and PostgreSQL instances.

\begin{lstlisting}[language=bash, caption={Run integration tests.}]
./gradlew integrationTest         # All integration tests
./gradlew integrationTest -i kafkaintegration   # Kafka only
\end{lstlisting}

\section{Performance Tests}

JMH benchmarks measure throughput and latency for critical paths:

\begin{lstlisting}[language=bash, caption={Run JMH benchmarks.}]
./gradlew jmh
\end{lstlisting}

%=============================================================================
\chapter{Pull Request Process}
\label{sec:pr-process}
%=============================================================================

\section{Branch Naming}

Use the following convention:

\begin{lstlisting}[language=bash, caption={Branch naming conventions.}]
feat/snowflake-sink-connector
fix/null-record-handling
docs/contributing-guide-update
\end{lstlisting}

\section{PR Checklist}

Before submitting a pull request, verify:

\begin{enumerate}
    \item All unit tests pass (\texttt{./gradlew test}).
    \item Code compiles without warnings (\texttt{./gradlew build}).
    \item New code has corresponding unit tests.
    \item Public APIs have Javadoc.
    \item Documentation is updated if behaviour changed.
    \item Commit messages follow Conventional Commits.
\end{enumerate}

\section{Review Process}

\begin{enumerate}
    \item Open the PR against \texttt{main}.
    \item Assign at least one reviewer from the module owners list.
    \item Address all review comments.
    \item CI must pass (build, test, lint, security scan).
    \item Squash-merge after approval.
\end{enumerate}

\section{Module Ownership}

\begin{table}[htbp]
    \centering
    \caption{Module owners and reviewers.}
    \label{tab:owners}
    \begin{tabular}{@{}lll@{}}
        \toprule
        Module & Owner & Reviewers \\
        \midrule
        core/ & @alex-feng & @mira-patel \\
        connectors/ & @jordan-lee & @alex-feng \\
        transforms/ & @sam-cho & @jordan-lee \\
        cli/ & @priya-nair & @sam-cho \\
        admin-api/ & @mira-patel & @priya-nair \\
        \bottomrule
    \end{tabular}
\end{table}

%=============================================================================
\chapter{Release Process}
\label{sec:release}
%=============================================================================

Releases follow semantic versioning. The release process is automated
via GitHub Actions:

\begin{enumerate}
    \item Update \texttt{CHANGELOG.md} with all changes since the last
        release.
    \item Create a release PR with the version bump in
        \texttt{build.gradle}.
    \item After approval, merge the PR.
    \item GitHub Actions creates a Git tag, builds the release artifacts,
        and publishes to Maven Central.
    \item Docker images are pushed to \texttt{ghcr.io/datapipe/datapipe}.
\end{enumerate}

\printbibliography[heading=bibintoc, title={References}]

\end{document}
